\documentclass{my_memo}

\usepackage{booktabs}
\usepackage{array}
\newcolumntype{L}[1]{>{\raggedright\arraybackslash}p{#1}}
\usepackage{hyperref}
\usepackage{xcolor}
\usepackage{listings}
\sloppy

\lstset{
  language=Fortran,
  basicstyle=\ttfamily\small,
  keywordstyle=\color{blue!60!black},
  commentstyle=\color{green!40!black},
  stringstyle=\color{red!60!black},
  columns=fullflexible,
  breaklines=true,
  frame=single,
  framesep=4pt,
  showstringspaces=false,
  numbers=none
}
\lstdefinestyle{sh}{language=bash,keywordstyle=\color{black},basicstyle=\ttfamily\small}

\newcommand{\code}[1]{\texttt{#1}}
\newcommand{\lamIlam}{\ensuremath{\lambda I_\lambda}}
\newcommand{\Nh}{\ensuremath{N_{\rm H}}}
\newcommand{\Cabs}{\ensuremath{C_{\rm abs}}}
\newcommand{\Teq}{\ensuremath{T_{\rm eq}}}
\newcommand{\um}{\ensuremath{\,\mu\mathrm{m}}}

\title{SEDust: User Manual}
\author{Kwang-il Seon}
\date{2026 August 1}

\begin{document}
\maketitle

\begin{abstract}
\noindent
This manual describes the model-agnostic dust library under \code{SEDust/sed/}.
It computes the infrared emission spectrum \emph{and} the transport optics
(extinction, scattering, asymmetry) of a dust mixture for an arbitrary local
radiation field, and is intended to be linked into a Fortran 3D
radiative-transfer (RT) code. It handles several dust models as peers --- the
HD23 \emph{astrodust}+PAH model, the Draine \& Li (2007) silicate+carbonaceous
model, the Zubko, Dwek \& Arendt (2004) BARE-GR-S model, and arbitrary
file-defined models --- through one object type, one emission call, and one
extinction call. The numerically validated stochastic-heating solver is reused
unchanged; the library is an additive layer on top of it.

This version is the \textbf{polarized} branch (v1.20). It is a
\emph{superset} of the scalar branch (v1.00): every scalar capability
documented here is present there and computes the same numbers, and v1.20 adds
the polarized cross sections, the direction-dependent extinction matrix and the
aligned-grain scattering matrix (\S\ref{sec:pol}). The one non-additive
difference is that the \code{status} codes of \code{sed\_init} /
\code{build\_astrodust} are numbered differently, because codes 3--5 are taken
by the polarized inputs: what is \code{3} and \code{4} in v1.00 is \code{6} and
\code{7} here (\S\ref{sec:status}). Code that only tests \code{status /= 0} is
unaffected.
\end{abstract}

\tableofcontents

%====================================================================
\section{Overview}

A \emph{dust model} is represented by a single derived type,
\code{dust\_model\_t}. You build it once (per RT run), then call
\code{dust\_emission} once per RT cell with that cell's local mean intensity:

\begin{lstlisting}
use dust_lib
type(dust_model_t)    :: m
real(wp), allocatable :: J_lam(:), lamI_total(:)

call build_astrodust(m, qtab, sizedist, 200, 2.7_wp, 5.0e3_wp)   ! once
allocate(J_lam(m%NLAM), lamI_total(m%NLAM))
do icell = 1, ncells
   ! ... assemble J_lam(:) (mean intensity on the grid m%lam) for this cell ...
   call dust_emission(m, J_lam, lamI_total)                      ! per cell
   ! ... lamI_total(:) is lambda*I_lambda / N_H for the cell ...
end do
\end{lstlisting}

\paragraph{Design.} The library wraps the existing, FIR-validated solver core
(\code{sed\_grain\_loop}, \code{calc\_Teq}, \code{calc\_P}) without modifying it.
A model is a set of grain \emph{populations} (one per species and charge state),
each carrying its own size distribution, absorption cross sections, enthalpy, and
Planck-integral tables. \code{dust\_emission} loops the populations through the
solver and sums their emission into named output \emph{channels}. The same
populations carry the extinction-side optics, so \code{dust\_extinction}
(\S\ref{sec:ext}) returns the model's opacity from that one object as well, and
emission and extinction both reach the transfer code through calls on one
model. Its scalar cross sections are served out of a size-integrated table the
builder attached (\S\ref{sec:ext}); the polarized ones are computed on every
call, because they follow the host's runtime alignment state.

\paragraph{One active model at a time.} The module-global grids
(\code{lam}, \code{T\_first}, $\dots$) hold the \emph{currently built} model's
working set. Calling \code{build\_*} makes that model active. This is sufficient
for an RT run that uses a single dust model; to switch models, rebuild.

%====================================================================
\section{Building the library}

From \code{SEDust/sed/}:
\begin{lstlisting}[style=sh]
make libsedust.a      # the library archive (link this into your RT code)
\end{lstlisting}

The compiler is \code{gfortran}. The default flags are
\code{-O3 -ffree-line-length-none -fopenmp} plus a warning set --- deliberately
\emph{portable}, so that a \code{libsedust.a} built on one node runs on any
other. A \code{-march=native} line is present but commented out; enabling it can
change the last digits through FMA contraction. \code{libsedust.a} is a static
archive of all library objects; the Fortran \code{.mod} files are written
alongside it in \code{SEDust/sed/}. Because those \code{.mod} files share one
directory, the \code{Makefile} carries \code{.NOTPARALLEL} and builds serially;
\code{make -j} is not supported until the object and module graph lives in a
separate directory.

%====================================================================
\section{Models and builders}

Each builder fills a \code{dust\_model\_t} and defines the model's output
channels (Table~\ref{tab:models}).

\begin{table}[h]
\centering
\begin{tabular}{@{}llll@{}}
\toprule
builder & model & channels & notes \\
\midrule
\code{build\_astrodust} & HD23 astrodust+PAH & \code{AD}, \code{PAH} & $N_C{=}417$, D16 graphite \\
\code{build\_dl07}      & Draine \& Li 2007   & \code{SIL}, \code{CARB} & $N_C{=}470$, D03 graphite \\
\code{build\_zubko}     & Zubko (ZDA 2004) BARE-GR-S & \code{PAH}, \code{GRA}, \code{SIL} & neutral PAH only \\
\code{build\_from\_files} & file-defined & \code{CH1}, \code{CH2}, $\dots$ & generic loader \\
\bottomrule
\end{tabular}
\caption{Built-in model builders. Each sets the model's own optics/abundance
parameters internally.}
\label{tab:models}
\end{table}

\noindent Signatures (all optional arguments in brackets; every optional is a
keyword argument in the order shown):
\begin{lstlisting}
call build_astrodust(m, qtable_path, sizedist_path, NT, T_lo, T_hi &
                     [, status] [, qpol_path] [, qpol_wave_path] [, qpol_aeff_path] &
                     [, scatmat_path] [, load_polarized_optics] &
                     [, lam_min] [, astrodust_index_path] &
                     [, qpol_euv_path] [, qpol_euv_wave_path])
call build_dl07     (m, qtable_path, sizedist_path, sd_index, U, NT, T_lo, T_hi &
                     [, status] [, lam_min])
call build_zubko    (m, config_path, data_dir, NT, T_lo, T_hi [, status])
call build_from_files(m, descriptor_path, data_dir, NT, T_lo, T_hi [, status])
\end{lstlisting}

\noindent The polarized arguments (\code{qpol\_*}, \code{scatmat\_path},
\code{load\_polarized\_optics}) are described in \S\ref{sec:pol}; the rest are
common to both branches.

\begin{itemize}
\item \code{NT}, \code{T\_lo}, \code{T\_hi}: number and range [K] of the internal
      temperature grid (typical \code{200, 2.7, 5000}).
\item \code{sd\_index} (DL07): WD01 size-distribution index --- \code{7} =
      MW $R_V{=}3.1$, $q_{\rm PAH}{=}4.58\%$; \code{29/30/31} = LMC2; \code{32} = SMC.
\item \code{U} (DL07): ISRF intensity scaling (the DL07 reference uses $U{=}1$).
\item \code{data\_dir} (Zubko / files): directory holding the data files, e.g.
      \code{data/zubko/}.
\item \code{status} (optional, all builders): error report; see
      \S\ref{sec:status}. When present, a missing or malformed input file is
      reported through it and the model is not built, so an RT host can recover.
\item \code{lam\_min} (optional, astrodust and DL07): shortest wavelength
      [\um] the model must cover, for a host that transports ionizing
      radiation. See \S\ref{sec:euv}. \textbf{Omitting it reproduces the previous
      behavior exactly} --- the grid, and every number derived from it, is
      unchanged.
\item \code{astrodust\_index\_path} (optional, astrodust only): the dielectric
      function that supplies the optics below $0.0912$\um. It must be the file
      \code{qtable\_path} was computed from --- same porosity, iron fraction and
      axial ratio --- or the model changes material at the seam. Omitted, it is
      the default that pairs with the shipped $Q$ table:
\begin{lstlisting}[style=sh]
data/dielectric/index_DH21Ad_P0.20_0.00_1.400   <->
tmatrix/output/q_astrodust_P0.20_Fe0.00_1.400.dat
\end{lstlisting}
\end{itemize}

\paragraph{Two routes to the Zubko model.} \code{build\_zubko} evaluates the ZDA
log-polynomial size-distribution \emph{formula} (coefficients read from
\code{ZDA\_BARE\_GR\_S\_Config.dat}) on the optics-table radii;
\code{build\_from\_files} with \code{data/zubko/zubko\_descriptor.txt} instead
interpolates the \emph{tabulated} \code{SzDist} files onto the same radii. The
two shapes agree to $\sim10^{-5}$, but the distributed tables carry a
normalization that sits below $A\,g(a)$ of the config by 0.79\% (PAH), 1.13\%
(graphite) and 0.41\% (silicate), and they are sampled at only 24 / 62 / 63
radii, so interpolation near $a_{\max}$ adds a little more. Together this moves
the size-integrated $C_{\rm ext}$ by up to 2.10\% (median 1.02\%). \textbf{The
formula is the default and is faithful to the ZDA definition}; the tabulated
route exists to reproduce the benchmark's own size grid.

%====================================================================
\section{Computing emission}

\subsection{Single-cell exact solve}

\begin{lstlisting}
call dust_emission(m, J_lam, lamI_total [, lamI_chan] [, status])
\end{lstlisting}
\begin{itemize}
\item \code{J\_lam(m\%NLAM)}: the local \emph{mean intensity} on the grid
      \code{m\%lam} [\um]. For a scaled Mathis ISRF use
      \code{call J\_Mathis(U, m\%lam, J\_lam)} from module \code{radfield}.
\item \code{lamI\_total(m\%NLAM)}: the summed SED, \lamIlam$/\Nh$
      [\,erg\,s$^{-1}$\,cm$^{-2}$\,sr$^{-1}$\,H$^{-1}$\,] --- the HD23 convention.
\item \code{lamI\_chan(m\%NLAM, m\%n\_channel)} (optional): the SED of each channel.
\item \code{status} (optional): error report; see \S\ref{sec:status}.
\end{itemize}

\subsection{Solver choice (\code{m\%stoch\_method})}

The field \code{m\%stoch\_method} selects how each grain is heated
(Table~\ref{tab:solver}); \code{dust\_emission} honors it.

\begin{table}[h]
\centering
\begin{tabular}{@{}lll@{}}
\toprule
\code{m\%stoch\_method} & method & note \\
\midrule
\code{'heuristic'} & GD look-ahead T-window narrowing & \textbf{default}; fastest, serial \\
\code{'draine'}    & Guhathakurta \& Draine (1989) & reference GD \\
\code{'qm'}        & energy-space transition matrix (thermal-discrete) & most detailed \\
\code{'equil'}     & single-grain equilibrium (no stochastic) & see \S\ref{sec:eqtemp} \\
\bottomrule
\end{tabular}
\caption{Stochastic-heating solver options. The first three resolve the grain
temperature \emph{distribution} (PAH/NIR features); \code{'equil'} forces
equilibrium.}
\label{tab:solver}
\end{table}

\subsection{Equilibrium-temperature options}\label{sec:eqtemp}

Two options replace the full stochastic solve when an equilibrium approximation
is wanted:

\paragraph{(1) Equilibrium temperature for each grain type and size.}
Set \code{m\%stoch\_method = 'equil'} and call \code{dust\_emission}. Every grain
is placed at its own equilibrium temperature \Teq, computed from \emph{that
grain's} absorption cross section, and emits $C_{\rm abs}\,B_\lambda(\Teq)$. No
stochastic excursions are computed.

\paragraph{(2) A single equilibrium temperature for the whole model.}
\begin{lstlisting}
call dust_emission_single_teq(m, J_lam, lamI_total [, Teq_out])
\end{lstlisting}
All dust shares one temperature, found from the type- and size-integrated
(\,dn-weighted\,) total absorption cross section
$C_{\rm abs}^{\rm tot}(\lambda) = \sum_{\rm pop}\sum_a {\rm d}n(a)\,C_{\rm abs}(\lambda,a)$.
The single \Teq\ satisfies $\int C_{\rm abs}^{\rm tot}\,J\,{\rm d}\lambda =
\int C_{\rm abs}^{\rm tot}\,B(\Teq)\,{\rm d}\lambda$, and the emission is
$C_{\rm abs}^{\rm tot}\,B_\lambda(\Teq)$. The optional \code{Teq\_out} returns that
temperature.

Both options conserve energy (the emitted bolometric power equals the absorbed
power); they differ from the stochastic solve only in the SED \emph{shape} ---
a smooth modified blackbody rather than one carrying PAH/NIR features.

\subsection{Precomputed table + interpolation}

When the field in each cell is a fixed spectral \emph{shape} scaled by an intensity
$U$, precompute a table once and interpolate per cell:
\begin{lstlisting}
type(dust_emis_table_t) :: tab
call dust_build_table(m, J_ref, U_grid, tab [, status])            ! once
call dust_emission_interp(tab, U, lamI_total [, lamI_chan] [, status])  ! per cell
\end{lstlisting}
\code{J\_ref(m\%NLAM)} is the reference field shape (at $U{=}1$) and
\code{U\_grid(:)} an ascending intensity grid. The table is built from exact
solves, so it is exact at the grid points; interpolation is log--log in $U$.
The stored emissivities are floored at $10^{-300}$ before the log, so a grid
node whose true value is $0$ comes back as $10^{-300}$ rather than exactly $0$.
Both calls take an optional final \code{status} (\S\ref{sec:status}).

\paragraph{Caveat.} Stochastic emission is a \emph{nonlinear} functional of the
full $J(\lambda)$, not of a scalar. The table is valid only when the cell field
is $U\times(\text{fixed }J_{\rm ref})$. For cells whose field \emph{shape}
departs from $J_{\rm ref}$ (e.g.\ hardened spectra near hot stars), use the
cell-by-cell exact \code{dust\_emission}.

\subsection{Optional error status}\label{sec:status}

The builders and the emission/table calls each accept an optional final
\code{status} argument (\code{integer, intent(out)}). When it is present, an
invalid model, argument, or input file is reported through it (\code{status = 0}
on success, nonzero on failure) and the process keeps running, so a 3D RT host
can recover; when it is omitted, the same condition prints a message and
\code{stop}s the run, which the CLI drivers rely on. For the builders the
nonzero value only names the failing stage: the contract is simply
\code{0} = built, nonzero = build failed.

\begin{itemize}
\item \code{dust\_emission}: \code{1} unknown \code{stoch\_method};
      \code{2} \code{'qm'} selected but a population's radii are unset.
\item \code{dust\_build\_table}: \code{1} \code{size(J\_ref)}$\,\ne$\code{m\%NLAM};
      \code{2} \code{size(U\_grid)}$\,<2$;
      \code{3} \code{U\_grid} not positive and strictly increasing.
\item \code{dust\_emission\_interp}: \code{1} \code{U}$\,\le0$;
      \code{2} \code{size(lamI\_total)}$\,\ne$\code{tab\%NLAM};
      \code{3} \code{lamI\_chan} present but not
      \code{(tab\%NLAM, tab\%n\_channel)}.
\item \code{build\_dl07}: \code{1} $Q$-table load failed;
      \code{2} size-distribution load failed;
      \code{7} \code{lam\_min} is shorter than the D03 dielectric functions
      ($6.1992\times10^{-5}$\um\ --- the longer-reaching of astrosilicate and
      graphite), which is refused rather than served with a refractive index
      frozen at the table boundary. Codes \code{3}--\code{6}, \code{8} and
      \code{9} cannot occur on this path: the DL07 model has no polarized
      optics and needs no separate EUV dielectric function, because its optics
      are D03 Mie at every wavelength. The numbering is shared with
      \code{build\_astrodust} on purpose, so a host can decode one table.
\item \code{build\_astrodust} (and \code{sed\_init}):
      \code{1} $Q$-table load failed;
      \code{2} size-distribution load failed;
      \code{3} an explicit \code{scatmat\_path} failed to load;
      \code{4} an explicit \code{qpol\_path} could not be read (an implicit
      default degrades gracefully instead);
      \code{5} \code{load\_polarized\_optics=.false.} combined with an explicit
      \code{qpol\_*}/\code{scatmat\_path} (a contradiction);
      \code{6} the astrodust dielectric function failed to load (raised only
      when \code{lam\_min} asks for the EUV band);
      \code{7} \code{lam\_min} is shorter than that dielectric function's own
      shortest wavelength ($1.00003\times10^{-4}$\um), refused for the same
      reason as above;
      \code{8} the EUV companion polarized table could not be read while the
      grid reaches into that band and polarized optics were asked for;
      \code{9} the EUV part of the grid runs outside the wavelengths that table
      covers ($0.0124$--$0.0912$\um).
      \textbf{Codes 6 and 7 are 3 and 4 in the scalar branch (v1.00)}, which
      has no polarized inputs to occupy 3--5.
\item \code{build\_zubko}: \code{1} config read; \code{2} fewer than 3 components;
      \code{3} a component's optics read; \code{4} grid inconsistent;
      \code{5} a component's calorimetry read failed;
      \code{6} an explicitly named \code{kext\_path} failed to load.
\item \code{build\_from\_files}: \code{1} descriptor open;
      \code{2} too many \code{pop:} lines; \code{3} invalid channel;
      \code{4} no \code{pop:} lines; \code{5} optics read;
      \code{6} grid inconsistent; \code{7} size-distribution read;
      \code{8} calorimetry read failed;
      \code{9} an explicitly named \code{kext\_path} failed to load.
\item \textbf{An explicitly named \code{kext\_path} that fails to load}
      (\S\ref{sec:ext}) is \code{5} in \code{build\_dl07}, \code{6} in
      \code{build\_zubko}, \code{9} in \code{build\_from\_files} --- the next
      free code in each --- and \code{10} in \code{build\_astrodust}, whose
      polarized codes already occupy 3, 4, 5, 8 and 9. The scalar branch
      (v1.00) uses the same numbers except that its \code{build\_astrodust}
      has no polarized codes and so uses \code{5} there too.
\item \code{dust\_extinction}: \code{1} an output array is not of size
      \code{m\%NLAM}; \code{2} no extinction table was loaded for this model;
      \code{3} \code{m\%lam} runs outside the table.
\end{itemize}

%====================================================================
\section{Computing extinction}\label{sec:ext}

\code{dust\_extinction} is the extinction counterpart of \code{dust\_emission}:
a transfer code takes its opacity from the same model object, and on the same
wavelength grid \code{m\%lam}, as its emission.

\begin{lstlisting}
call dust_extinction(m, Cext, Cabs, Csca [, gbar] [, Cpol_ext] [, Cbir_ext] &
                     [, albedo] [, status])
\end{lstlisting}
\begin{itemize}
\item \code{Cext}, \code{Cabs}, \code{Csca} \code{(m\%NLAM)}: extinction,
      absorption and scattering cross sections per H atom [\,cm$^2$/H\,], with
      $C_{\rm ext} = C_{\rm abs} + C_{\rm sca}$.
\item \code{gbar(m\%NLAM)} (optional): the scattering asymmetry
      $\langle\cos\theta\rangle$.
\item \code{Cpol\_ext(m\%NLAM)} (optional): the dichroic (polarized)
      extinction [\,cm$^2$/H\,]; see \S\ref{sec:pol}.
\item \code{Cbir\_ext(m\%NLAM)} (optional): the birefringent extinction
      [\,cm$^2$/H\,], the $f_{\rm align}$-weighted size integral of the
      orientation-resolved table's 4th block. It is $0$ when the loaded table
      carries no birefringence (the release table does not); see \S\ref{sec:pol}.
\item \code{albedo(m\%NLAM)} (optional): $C_{\rm sca}/C_{\rm ext}$, and $0$
      where $C_{\rm ext}$ underflows to zero. It is derived here rather than by
      the caller so that every host gets the same convention at those
      wavelengths.
\item \code{status} (optional): error report; see \S\ref{sec:status}.
      It is $0$ on success, $1$ if an output array is not of size
      \code{m\%NLAM}, $2$ if no extinction table was loaded for this model, and
      $3$ if \code{m\%lam} runs outside the table.
\end{itemize}

\paragraph{Where the numbers come from.} The four \emph{scalar} outputs ---
\code{Cext}, \code{Cabs}, \code{Csca}, \code{gbar} --- are read from a
precomputed size-integrated extinction table, one of the
\code{data/kext\_*.dat} products of \code{calc\_kext.x}
(\S\ref{sec:standalone}), attached to the model at build time and interpolated
onto \code{m\%lam}. The transport optics of a dust model do not depend on the
transport, so there is nothing to gain by repeating the size integral during a
transfer run: the table \emph{is} that integral, done once and recorded.

The two \emph{polarized} outputs --- \code{Cpol\_ext} and \code{Cbir\_ext} ---
are \textbf{not} tabulated. They are computed from the model's own
orientation-resolved optics on every call, because their $f_{\rm align}(a)$
weight is runtime state a transfer code resets cell by cell through
\code{dust\_set\_alignment} (\S\ref{sec:pol:align}), which a table fixed at
build time could not follow. This asymmetry --- scalar optics from a table,
polarized optics from the model --- is deliberate.

The first-principles size integral remains available under its own name and
with an identical argument list:
\begin{lstlisting}
call size_integrated_extinction(m, Cext, Cabs, Csca [, gbar] [, Cpol_ext] &
                                [, Cbir_ext] [, albedo] [, status])
\end{lstlisting}
It needs no table and computes all six outputs from the model's optics. It is
what \code{calc\_kext.x} calls to write the tables in the first place, and what
the in-tree checks (\code{test\_scalar\_mode.x},
\code{test\_euv\_extension.x}) call when they have to compare one build of a
model against another. Each output of it is the plain binned sum over every
population of the model, since ${\rm d}n(a)$ already carries the bin width:
\begin{align}
C_{\rm abs}(\lambda) &= \sum_{\rm pop}\sum_a {\rm d}n(a)\,C_{\rm abs}(\lambda,a), &
C_{\rm sca}(\lambda) &= \sum_{\rm pop}\sum_a {\rm d}n(a)\,C_{\rm sca}(\lambda,a), \nonumber\\
\langle\cos\theta\rangle &= \frac{\sum {\rm d}n\,C_{\rm sca}\,g}{\sum {\rm d}n\,C_{\rm sca}}, &
C_{\rm pol,ext}(\lambda) &= \sum_{\rm pop}\sum_a {\rm d}n(a)\,C_{\rm pol,ext}(\lambda,a)\,f_{\rm align}(a).
\label{eq:extsums}
\end{align}
A population that carries no scattering or polarized optics contributes zero to
those terms. In the astrodust model the PAHs are exactly that case: they enter
through absorption only, and the astrodust grains carry all the scattering and
the asymmetry.

\paragraph{Dust mass per H, and mass opacity.} Both routines
return cross sections \emph{per H nucleon}. A host that carries a dust mass
density instead of a hydrogen column converts them with one number, which the
model itself reports:
\begin{lstlisting}
Mdust_H = dust_mass_per_H(m)   ! [g/H]
kappa   = Cabs / Mdust_H       ! [cm^2 per gram of dust]
\end{lstlisting}
It is the model's own size distribution weighted by the solid density of each
population's material,
\begin{equation}
\frac{M_{\rm dust}}{N_{\rm H}}
  \;=\; \sum_{\rm pop} \rho_{\rm bulk}
        \sum_a \tfrac43\pi a^3\,{\rm d}n_{\rm pop}(a),
\label{eq:massperh}
\end{equation}
with $a$ in cm and ${\rm d}n$ already carrying the bin width, so the opacity it
produces refers to the same grains the emission does. Being a property of the
model it is independent of wavelength and temperature: evaluate it once. A
population whose builder states no density ($\rho_{\rm bulk}=0$) is skipped
rather than given a guessed one, and a model in which no population states a
density returns $0$.

The densities are part of each model's definition, not a free choice --- the
size distributions were normalized to element abundances with them:

\begin{center}
\begin{tabular}{@{}lll@{}}
\toprule
model & population & $\rho_{\rm bulk}$ [g\,cm$^{-3}$] \\
\midrule
astrodust  & astrodust grains & $2.74$ (HD23 \S2.3.2, already $\times(1-P)$) \\
           & PAH              & $2.0$ (HD23 $N_{\rm C}=417\,(a/{\rm nm})^3$) \\
DL07       & astrosilicate    & $3.5$ (DL07 \S2) \\
           & carbonaceous     & $2.2$ (graphitic carbon, DL07 \S2) \\
Zubko, file-defined & each component & declared in its own optics table \\
\bottomrule
\end{tabular}
\end{center}

\noindent Both DL07 values are the ones its own paper states, in the \S2
paragraph that defines the effective radius $a = (3M/4\pi\rho)^{1/3}$:
``amorphous silicate is assumed to have a mass density $\rho = 3.5$\,g\,cm$^{-3}$,
and carbonaceous grains are assumed to have a mass density due to graphitic
carbon alone of $\rho = 2.2$\,g\,cm$^{-3}$''. The graphite value is
deliberately \emph{not} the $2.24$ of WD01/LD01, which is the density implied by
the carbon-atom count $N_{\rm C} = 470\,(a/{\rm nm})^3$ this builder selects;
DL07 kept that count and quoted the rounder density, and the model built here is
DL07's.

The DL07 carbonaceous grains are one continuous material sequence --- PAH-like
when small, graphite-like when large, joined by the DL07 $\xi(a)$ weight ---
with no size at which the solid changes, so the graphite density is used across
the whole sequence rather than split at some radius. That the astrodust model's
PAHs use $2.0$ instead is HD23's convention for a different material in a
different model, not a disagreement about graphite.

The same number normalizes the $K_{\rm abs}$ column of every
\code{data/kext\_*.dat} product (\S\ref{sec:standalone}), so a host that reads
the tables and a host that links the library get the identical conversion. The
values are $1.184949\times10^{-26}$ g/H for astrodust,
$1.750186\times10^{-26}$ g/H for DL07 and $1.443932\times10^{-26}$ g/H for
Zubko BARE-GR-S.

\paragraph{Which DL07 normalization is followed, and why it matters.} DL07's
own Table~3 gives $M_{\rm dust}/M_{\rm H} = 0.0104$ for model $j_M = 7$
(MW3.1\_60) --- the very size distribution \code{build\_dl07} selects --- and the
densities above reproduce it to about half a percent
($1.750186\times10^{-26}$ g/H, i.e.\ $M_{\rm dust}/M_{\rm H} = 0.01046$ with
$m_{\rm H} = 1.6737\times10^{-24}$ g, $+0.55$\%). Draine's tabulated optics
for the same model, \code{kext\_albedo\_WD\_MW\_3.1\_60\_D03.all\_2003}, state
$M_{\rm dust}/N_{\rm H} = 1.870\times10^{-26}$ g/H in the file header, i.e.
$M_{\rm dust}/M_{\rm H} = 0.0112$, which is 7\% away from that author's own
Table~3. The same file renormalizes the silicate distribution by $0.93$ where
the paper's Fig.~11 caption says $0.92$.

\noindent SEDust takes the \emph{optics} from the file --- reproducing it is the
point, and our $C_{\rm ext}$/H agrees with it to $0.1$--$0.3$\% over its whole
range --- and the \emph{mass} from the paper. The two normalizations differ, so
about 1\% remains between our $K_{\rm abs}$ column and the file's; that residual
is on Draine's side, not a free parameter here.

\paragraph{Choosing the table.}\label{par:kextpath} Every builder takes an
optional \code{kext\_path} naming the file. Omitting it takes that model's
default:

\begin{center}
\begin{tabular}{@{}ll@{}}
\toprule
builder & default table \\
\midrule
\code{build\_astrodust}  & \code{../data/kext\_astrodust\_MW\_euv.dat} \\
\code{build\_dl07}       & \code{../data/kext\_dl07\_MW\_euv.dat} \\
\code{build\_zubko}      & \code{../data/kext\_zubko\_BARE\_GR\_S.dat} \\
\code{build\_from\_files}& none \\
\bottomrule
\end{tabular}
\end{center}

\noindent The astrodust and DL07 defaults are the EUV products because they are
the \emph{widest} grid each model has ($\sim10^{-4}$--$39810$\um): one file then
serves a host that transports ionizing radiation and a host that does not,
since the latter's grid is a subset of the former's, on the same nodes. Naming
a \code{kext\_path} is therefore how you \emph{narrow} the table (to
\code{kext\_astrodust\_MW.dat} or \code{kext\_dl07\_MW.dat}), or point at a
product of your own. A file-defined model has no default at all: its product is
named after the model, which only the descriptor knows, so a host that wants
extinction out of such a model must name the file.

A \code{kext\_path} that cannot be read \emph{fails the build}
(\S\ref{sec:status} for the code, which differs by builder) --- a host naming a
file that is not there is a configuration error.
A \emph{default} that cannot be read does not fail the build: the
emission-only drivers must still run, and \code{calc\_kext.x} builds a model
precisely in order to write the table that is not there yet. In that case
\code{m\%kext\_n} is $0$ and \code{dust\_extinction} reports status \code{2}.

The table is read at \emph{build} time, not at query time, because these paths
are relative to \code{sed/}. A host that changes directory around the
\code{build\_*} call --- which is what MoCafe does --- would otherwise find them
broken by the time it asks for opacity.

\paragraph{Interpolation, and when it is exact.} \code{Cext}, \code{Cabs} and
\code{Csca} are interpolated linearly in $\log\lambda$--$\log C$;
\code{gbar} linearly in $\log\lambda$, because it is signed and changes sign in
the far infrared of some models. A wavelength of \code{m\%lam} that coincides
with a table node (within $10^{-6}$ relative, the tolerance that absorbs the
round trip of $\lambda$ through a decimal file) takes the tabulated value
\emph{unchanged}. Nothing is extrapolated: a grid running outside the table is
refused with status \code{3}, not served a frozen boundary value.

An unextended astrodust or DL07 model therefore gets its numbers back
bit-for-bit: its 1129 wavelengths are the T-matrix $Q$-table grid, which is
exactly the long-wavelength part of the default EUV table (verified: all four
scalar outputs bit-identical at all 1129 nodes; \code{albedo} differs by
$\le10^{-12}$ relative because it is recomputed as $C_{\rm sca}/C_{\rm ext}$
rather than read). A model built with \code{lam\_min} does \emph{not} share the
table's nodes below $0.0912$\um, since a different floor gives a different
log spacing, and there the interpolation costs up to $\sim$1\% in $C_{\rm ext}$
(measured $1.2\times10^{-2}$ at worst, near $0.023$\um, for
$\code{lam\_min}=0.0124$\um). A host that needs the ionizing band at full
fidelity should generate a table at its own \code{lam\_min}
(\code{./calc\_kext.x astrodust <lam\_min>}) and pass it as \code{kext\_path}.

\paragraph{Every model scatters.} All four builders populate the
extinction-side scattering optics, so \code{albedo} and \code{gbar} are physical
for every model:

\begin{itemize}
\item \textbf{astrodust}: $C_{\rm sca}$ and $\langle\cos\theta\rangle$ from the
      random-orientation T-matrix $Q$ table (and, below $0.0912$\um, from Mie on
      the same dielectric function).
\item \textbf{DL07}: Mie on the D03 astrosilicate and graphite dielectric
      functions (\code{q\_silicate\_full} / \code{q\_graphite\_full}), on the
      same random-orientation average ($\tfrac13\,E\!\parallel\!c +
      \tfrac23\,E\!\perp\!c$) as the absorption. The two carbonaceous charge
      states share one scattering description: charge shifts PAH absorption
      features, not the graphite sphere's scattering.
\item \textbf{Zubko / file-defined}: the $Q_{\rm sca}$ and $g$ columns of the
      model's own DustEM tables, i.e.\ the authors' own multilayer-sphere
      solution, read from the same file as $Q_{\rm abs}$.
\end{itemize}

\noindent A population that genuinely does not scatter leaves its scattering
optics unallocated and contributes zero. In the astrodust model the PAHs are
exactly that case --- PAH scattering is Rayleigh-negligible --- so they enter
extinction through absorption only, and the astrodust grains carry all the
scattering and the asymmetry.

\paragraph{Polarized optics are astrodust only.} \code{Cpol\_ext} and
\code{Cbir\_ext} come back as zeros from \code{build\_dl07},
\code{build\_zubko} and \code{build\_from\_files}, which carry no polarized
optics at all, and from an astrodust model built with
\code{load\_polarized\_optics = .false.}
\code{dust\_has\_polarized\_optics(m)} distinguishes the two cases.

%====================================================================
\section{Polarized dust optics}\label{sec:pol}

The library exports two polarized quantities for the HD23 astrodust model: a
polarized \emph{emission} spectrum through \code{dust\_emission}, and a
polarized \emph{extinction} cross section through \code{dust\_extinction} or,
equivalently, through the precomputed extinction table. Both
are \emph{intrinsic}: the size integral and the alignment weighting are done,
but the geometry of the local magnetic field is not applied. What remains for
the calling code is set out in \S\ref{sec:pol:division}.

This section documents the interface only. For the alignment physics, the
polarized transfer equation, and the sign and index conventions behind these
quantities, see the companion report \code{aligned\_grain\_polarization.pdf}
(\S3 for the optical properties, \S4 for alignment, \S5 for transfer).

\paragraph{Wavelength coverage.} \code{Cpol}, \code{Cpol\_ext} and
\code{Cbir\_ext} come from an orientation-resolved table that covers all 1129
wavelengths of the $Q$-table grid, $0.0912$--$39810$\um, so nothing has to be
regenerated to use them anywhere in that range. The \emph{default} table is the
HD23 release file \code{data/dielectric/q\_DH21Ad\_P0.20\_Fe0.00\_1.400.dat.gz};
SEDust's own regenerated table,
\code{tmatrix/output/q\_astrodust\_jori\_P0.20\_Fe0.00\_1.400.dat.gz}, is opt-in
through \code{qpol\_path} and is the one that carries the 4th (birefringence)
block, so \code{Cbir\_ext} is zero with the default and nonzero with it
(\S\ref{sec:pol:source}). The angle-resolved \emph{scattering} matrix is a
separate, coarser product --- five optical bands (\S\ref{sec:pol:scatmat}) ---
and shortward of $0.0912$\um\ the polarized extinction has an announced hole
(\S\ref{sec:euv-pol}).

\subsection{Polarized emission}

\begin{lstlisting}
call dust_emission(m, J_lam, lamI_total [, lamI_chan] [, status] [, lamI_pol])
\end{lstlisting}
\begin{itemize}
\item \code{lamI\_pol(m\%NLAM)} (optional): the intrinsic polarized emission,
      with the same shape and the same units as \code{lamI\_total}. It is the
      emission a population of perfectly aligned grains would radiate with
      their symmetry axes in the plane of the sky.
\item Only populations that carry both polarized optics and an alignment
      efficiency contribute. PAHs are taken to be unaligned by HD23 and add
      zero, as does any model built without polarized optics.
\item The argument is optional and comes last, so every existing call site
      remains valid.
\end{itemize}

\noindent The polarized channel is accumulated at every point where the total
emission is accumulated, using the same grain temperature weights, so a
stochastically heated grain contributes to both consistently under all four
values of \code{m\%stoch\_method}.

\subsection{Polarized extinction}

Two routes give the same quantity. The direct one is the optional
\code{Cpol\_ext} argument of \code{dust\_extinction} (\S\ref{sec:ext}),
\begin{lstlisting}
call dust_extinction(m, Cext, Cabs, Csca, gbar, Cpol_ext)
\end{lstlisting}
which returns
$\sum_a {\rm d}n_{\rm Ad}(a)\,C_{\rm pol,ext}(\lambda,a)\,f_{\rm align}(a)$
[\,cm$^2$/H\,] on \code{m\%lam}. Prefer it in a code that already holds a model
object. Unlike the four scalar outputs of that call, this one is computed from
the model's optics rather than read from a table, so it follows the alignment
efficiency the host has set (\S\ref{sec:pol:align}) --- a tabulated column could
not.

The second route is the file \code{data/kext\_astrodust\_MW.dat}, written by
\code{calc\_kext.x astrodust}, see \S\ref{sec:standalone}, which carries the
dichroic extinction as an eighth column (Table~\ref{tab:kextcols}). It is the
right choice for a tool that wants the tabulated optics without linking the
library, and it is valid only for the HD23 alignment fit it was written under.
Codes that read only the first seven columns are unaffected by its presence, and
\code{dust\_extinction} is one of them: it reads no polarized column from any
table.

\begin{table}[h]
\centering
\begin{tabular}{@{}clll@{}}
\toprule
col. & quantity & unit & note \\
\midrule
1 & \code{lambda}        & \um\        & wavelength grid \\
2 & \code{albedo}        &             & $C_{\rm sca}/C_{\rm ext}$ \\
3 & \code{<cos>}         &             & scattering asymmetry $g$ \\
4 & \code{C\_ext/H}      & cm$^2$/H    & total extinction \\
5 & \code{C\_abs/H}      & cm$^2$/H    & total absorption \\
6 & \code{C\_sca/H}      & cm$^2$/H    & total scattering \\
7 & \code{K\_abs}        & cm$^2$/g    & mass absorption coefficient \\
8 & \code{C\_polext/H}   & cm$^2$/H    & dichroic extinction (new) \\
\bottomrule
\end{tabular}
\caption{Columns of \code{kext\_astrodust\_MW.dat}. Columns 1--7 are written for
every model; column 7 is $C_{\rm abs}/{\rm H}$ divided by that model's dust mass
per H, Eq.~(\ref{eq:massperh}). Column 8 is
$\sum_a {\rm d}n_{\rm Ad}(a)\,C_{\rm pol,ext}(\lambda,a)\,f_{\rm align}(a)$,
the maximum polarized extinction, quoted before any geometric factor, and this
file is the only product that carries it.}
\label{tab:kextcols}
\end{table}

\subsection{Changing the alignment efficiency}\label{sec:pol:align}

By default the alignment efficiency is the HD23 fit. Two setters, both
re-exported by \code{dust\_lib}, replace it on a model that already exists.

\begin{lstlisting}
call dust_set_alignment(m, f_max, a_align, alpha_align [, status])
call dust_set_alignment_profile(m, aeff_in, falign_in [, status])
\end{lstlisting}

The first installs the HD23 power-law rolloff
\begin{equation}
f_{\rm align}(a) \;=\; \frac{f_{\max}}{1 + (a_{\rm align}/a)^{\alpha_{\rm align}}},
\label{eq:falignset}
\end{equation}
with $a_{\rm align}$ in \um. The second installs an arbitrary tabulated
profile: the caller supplies $(\code{aeff\_in}\,[\um],\ \code{falign\_in})$
pairs and each population's efficiency is interpolated from them onto its own
radius grid, linearly in $\log a$ and clamped at the ends of the table. This
is the route for prescriptions Eq.~(\ref{eq:falignset}) cannot express: a
RAT-derived Rayleigh reduction factor $R(a)$, or the GRADE-POL exponential
$f_{\max}[1 - \exp(-(0.5a/a_{\rm align})^3)]$.

Both refill \code{falign} only for populations that carry polarized optics.
PAHs, which HD23 take to be unaligned, are left untouched and keep
contributing zero.

\paragraph{No temperature solve is repeated.} This is the property worth
knowing before designing a parameter study. The alignment efficiency enters as
a weight in the size integral, \emph{outside} the temperature solution: it
multiplies the polarized accumulator and never enters $P(T)$. Changing the
alignment therefore does not invalidate any temperature distribution, and
\code{lamI\_total} is bit-for-bit unchanged across a setter call. A scan over
alignment parameters costs one \code{dust\_emission} call per point, not one
model build.

\paragraph{Negative efficiencies are allowed on purpose.}
\code{dust\_set\_alignment\_profile} requires \code{falign\_in} to lie in
$[-1, 1]$ and \emph{rejects} rather than clamps a value outside it, since
$|f| \le 1$ is the physical bound on a Rayleigh reduction factor and a value
beyond it is a caller error that silent clamping would hide. Negative values
inside the range are accepted: a grain with slow internal relaxation has
$Q_X = -0.1$, a negative reduction factor that flips the polarization
direction by $90^\circ$. That is the accepted origin of polarization parallel
to $\mathbf{B}$ at millimeter wavelengths and perpendicular to it in the
submillimeter, and clamping at zero would silently delete it.
\code{dust\_set\_alignment} keeps $f_{\max}$ in $[0, 1]$, the power law having
no such sign freedom.

With \code{status} present a bad argument sets it nonzero and returns; without
it, the run stops. The codes are 1, 2, 3 in the order the arguments are listed
above for \code{dust\_set\_alignment} ($a_{\rm align} \le 0$,
$\alpha_{\rm align} \le 0$, $f_{\max}$ outside $[0,1]$), and for
\code{dust\_set\_alignment\_profile} 1 for mismatched or too-short arrays,
2 for an \code{aeff\_in} that is not positive and strictly increasing, and
3 for a \code{falign\_in} value outside $[-1, 1]$.

\paragraph{Status code 4: a loaded aligned scattering table is now stale.}
If an aligned scattering table has been loaded (\S\ref{sec:pol:scatmat}) and the
requested profile differs from the one that table was integrated under, either
setter returns the new code \code{4}. This is \emph{non-fatal}: the
$f_{\rm align}$ update completes and the emission and \code{Cpol\_ext} channels
re-integrate live under the new profile, so those outputs always honor the latest
alignment. The code only signals that the loaded scattering matrices no longer
correspond to it, because they were size-integrated once under the recorded
profile. The sanctioned runtime variation of the aligned scattering optics is the
scalar $\eta$ (\S\ref{sec:pol:scatmat}); a genuinely different profile requires
regenerating the table with \code{run\_scatmat\_aligned.x profile=FILE}. A
tabulated profile never matches the recorded analytic one and so always raises
the flag.

\subsection{Where the polarized optics come from}\label{sec:pol:source}

The orientation-resolved spheroid cross sections behind \code{lamI\_pol} and
\code{Cpol\_ext} default to the Draine \& Hensley (2021) release table. They can
instead be taken from a table SEDust regenerates from the astrodust dielectric
function with its own T-matrix engine: \code{build\_astrodust} and \code{sed\_init}
accept an optional \code{qpol\_path} (with \code{qpol\_wave\_path} and
\code{qpol\_aeff\_path} for its grid axes) that overrides the default file. The
regenerated table matches the release to a few parts in $10^4$ where grains carry
weight and is written in the release table's own format, so nothing else in the
call changes. Its
one deliberate benefit is that it supplies the polarized cross sections from the
same exact optics as the total intensity, removing the mixed averaging convention
the release forces; see \code{sedust\_polarization\_implementation.pdf} \S7 and
\code{tmatrix/run\_q\_jori.x}. Leaving \code{qpol\_path} unset keeps the release
table and the existing results, though an explicit \code{qpol\_path} that cannot
be read now fails the build with \code{status = 4} (and an unconditional stderr
message), whereas the implicit default still degrades gracefully to an unpolarized
model.

\paragraph{Scalar-only builds.} \code{build\_astrodust} and \code{sed\_init} take
an optional \code{load\_polarized\_optics}. Absent or \code{.true.} is the present
behavior. \code{.false.} never opens the orientation-resolved polarized $Q$ table:
\code{Cpol}, \code{Cpol\_ext}, \code{Cbir\_ext} and \code{falign} stay zero, and
the scalar \code{Cext}/\code{Cabs}/\code{Csca}/\code{gbar} and the total SED come
back bit-identical to a polarized build at zero alignment. Combining it with an
explicit \code{qpol\_*} or \code{scatmat\_path} is a contradiction and returns
\code{status = 5}. The query \code{dust\_has\_polarized\_optics(m)} reports whether
the built model carries polarized optics at all; the aligned-scattering table's
load state is separately visible as \code{scm\_loaded}.

\subsection{Aligned-grain scattering optics}\label{sec:pol:scatmat}

For the astrodust spheroid SEDust also supplies the angle-resolved scattering
optics of the aligned population, for a polarized RT code that transports a Stokes
vector: the fixed-orientation Mueller matrix, the $4\times4$ extinction matrix,
and the random-orientation remainder. These come from a table written by
\code{tmatrix/run\_scatmat\_aligned.x} (five optical bands ship with SEDust); how
it is built is documented in \code{sedust\_polarization\_implementation.pdf}
\S8. The API follows the library's own lifecycle --- initialize once, then query
along the photon path --- split into two strictly separated layers.

\paragraph{Loading (once per run, serial).} Pass the table path as an optional
\code{scatmat\_path} to \code{sed\_init} or \code{build\_astrodust}, exactly like
\code{qpol\_path}:
\begin{lstlisting}
call build_astrodust(m, qtab, sizedist, NT, T_lo, T_hi, status=st, &
                     qpol_path=QPOL, scatmat_path=SCATMAT)
\end{lstlisting}
The production table holds five optical bands --- $0.36$, $0.44$, $0.55$, $0.64$
and $0.79$\um, approximately $UBVRI$ --- on a grid of 19 incidence angles
$\theta_i$ $\times$ 181 scattering angles $\theta_s$ $\times$ 37 azimuths
$\phi$. Five bands is a deliberate choice, not a limitation of the format:
polarized transfer is almost never wanted at every wavelength, and what matters
is that any other wavelength \emph{can} be produced (\S\ref{sec:standalone:qjori}
--- note the caveats there, since the scattering table is the less flexible of
the two products). At run time the queries are a trilinear interpolation on that
grid, a few hundred floating-point operations per scattering event, and
\code{scatmat\_band} maps a requested wavelength onto a stored band once.

The table is parsed and its size integrals kept in memory (about $81.5$\,MB for the
production grid, growing linearly with the number of bands). A
\code{scatmat\_path} that fails to load is an error
(\code{status = 3} from the builder), because a host that asked for the table
depends on it. A gzipped table --- both this and the polarized $Q$ table ship
\code{.gz} --- is decompressed to a scratch copy under \code{TMPDIR} (else
\code{/tmp}) named with the process id (\code{scatmat\_aligned\_<pid>.dat},
\code{q\_jori\_<pid>.dat}) and deleted after the read, so concurrent MPI ranks
never collide and a read-only launch directory is never written. This layer is not
thread-safe; call it from serial code before the photon loop.

\paragraph{Path queries (per event, thread-safe).} Six routines, all re-exported
by \code{dust\_lib}, are pure reads of the loaded arrays: they write only their own
output arguments, do no I/O and no allocation, and are safe to call concurrently
from OpenMP photon threads. The whole cell dependence enters through two scalars
the host already holds --- the alignment scale $\eta$, and the incidence angle
$\theta_i = \arccos(\hat k\!\cdot\!\hat B)$ between the photon direction and the
local field, from one dot product.
\begin{lstlisting}
call scatmat_band(lambda_um, iband, exact)                       ! pick the band once
call extinction_matrix_aligned(iband, theta_i, eta, kmat)        ! 4x4  [um^2/H]
call mueller_matrix_total(iband, theta_i, theta_s, phi, eta, z)  ! 4x4  [um^2 sr^-1/H], absolute
call mueller_matrix_aligned(iband, theta_i, theta_s, phi, z)     ! 4x4  [um^2 sr^-1/H], aligned only, eta=1
call mueller_matrix_random(iband, big_theta, f_tot, f_ref)       ! 6-elt random matrices
call scattering_cross_sections(iband, theta_i, eta, csca_aligned, csca_unaligned)
\end{lstlisting}
\begin{itemize}
\item \code{scatmat\_band}: \code{lambda\_um} [\um] in; \code{iband} the nearest
      stored band; \code{exact} \code{.true.} if it matched to $10^{-3}$
      fractionally. The hot path then takes \code{iband} so no wavelength search
      runs each event.
\item \code{extinction\_matrix\_aligned}: \code{kmat(4,4)} [\,\um$^2$/H\,], the
      $\eta$-scaled aligned extinction matrix at incidence \code{theta\_i} [deg].
      Diagonal $\eta\,C_{\rm ext,al}$; $K(1,2)=K(2,1)=\eta\,C_{\rm pol,al}$
      (dichroism); $K(3,4)=\eta\,C_{\rm bir,al}$, $K(4,3)=-\eta\,C_{\rm bir,al}$
      (birefringence). \code{theta\_i} in $[0,180^\circ]$ is folded to
      $[0,90^\circ]$, $K$ being even under the equatorial reflection.
\item \code{mueller\_matrix\_total} (\textbf{recommended}): the \emph{absolute}
      combined phase matrix \code{z(4,4)} [\,\um$^2$\,sr$^{-1}$/H\,] at
      grain-frame geometry and alignment $\eta$, with both populations already
      assembled,
      \[
        z = \eta\,Z_{\rm al}
          + L(\pi-\sigma_2)\,F_{\rm unal}(\Theta)\,L(-\sigma_1)\,/\,(4\pi),
        \quad
        F_{\rm unal}=C_{\rm sca,tot}F_{\rm tot}-\eta\,C_{\rm sca,ref}F_{\rm ref},
      \]
      with $\cos\Theta=\cos\theta_i\cos\theta_s+\sin\theta_i\sin\theta_s\cos\phi$
      and the meridional rotation angles (closed forms, poles included)
      \begin{align*}
        \sigma_1 &= \mathrm{atan2}\!\left(\sin\theta_s\sin\phi,\;
                    \cos\theta_i\sin\theta_s\cos\phi-\sin\theta_i\cos\theta_s\right),\\
        \sigma_2 &= \mathrm{atan2}\!\left(-\sin\theta_i\sin\phi,\;
                    \cos\theta_i\sin\theta_s-\sin\theta_i\cos\theta_s\cos\phi\right).
      \end{align*}
      $L$ is the Stokes rotation. Prefer this call to reconstructing the mixture
      from the public arrays: it is the one place the normalization and the
      rotation signs are certified.
\item \code{mueller\_matrix\_aligned}: the aligned matrix alone, \code{z(4,4)}
      [\,\um$^2$\,sr$^{-1}$/H\,] at the reference alignment ($\eta=1$; scale it
      yourself). \code{theta\_i}, \code{theta\_s} in $[0,180^\circ]$, \code{phi}
      in $[0,360^\circ]$; the stored octants ($\theta_i\!\le\!90^\circ$,
      $\phi\!\le\!180^\circ$) and the rest are handled by interpolation plus the
      header's symmetry relations.
\item \code{mueller\_matrix\_random}: \code{f\_tot(6)}, \code{f\_ref(6)}, the two
      six-element ($F_{11},F_{22},F_{33},F_{44},F_{12},F_{34}$) random-orientation
      matrices at \code{big\_theta} [deg], $\alpha_1$-normalized as stored
      ($\tfrac12\int F_{11}\,{\rm d}\cos\Theta=1$, i.e.\ $\int F_{11}\,{\rm
      d}\Omega=4\pi$). The absolute differential matrix is $C_{\rm sca}F/(4\pi)$:
      restore \um$^2$\,sr$^{-1}$/H with
      \code{f\_tot*scm\_csca\_tot(iband)/(4*PI)} and
      \code{f\_ref*scm\_csca\_ref(iband)/(4*PI)}, the $1/(4\pi)$ turning the
      $4\pi$-normalized $F$ into the matrix whose element 11 integrates to
      $C_{\rm sca}$ over the sphere. (Omitting it overweights the remainder by
      $4\pi$ relative to $Z_{\rm al}$, whose stored matrix already integrates to
      $C_{\rm sca,al}$.)
\item \code{scattering\_cross\_sections}: \code{csca\_aligned} $=\eta\,
      C_{\rm sca,al}(\theta_i)$ and \code{csca\_unaligned}
      $=C_{\rm sca,tot}-\eta\,C_{\rm sca,ref}$ [\,\um$^2$/H\,], for the albedo and
      the population split. An optional fourth output
      \code{csca\_pol\_aligned} $=\eta\,C_{\rm sca,pol,al}(\theta_i)$ is the
      polarized scattering cross section $\int Z_{12}\,{\rm d}\Omega$ of the
      aligned matrix (Peest et al.\ 2023), which the polarization-dependent albedo
      $\Lambda=[C_{\rm sca}+C_{\rm sca,pol}\,Q/I]/[C_{\rm ext}+C_{\rm pol}\,Q/I]$
      needs; the random remainder contributes zero to it by azimuthal symmetry.
\end{itemize}

\paragraph{The $\eta$ contract, and unit note.} A cell's alignment enters only as
$\eta$: the aligned matrix and the aligned $K$ scale as $\eta$, and the unaligned
remainder scattering matrix is
$Z_{\rm unal}=[C_{\rm sca,tot}F_{\rm tot}-\eta\,C_{\rm sca,ref}F_{\rm ref}]/(4\pi)$
in absolute \um$^2$\,sr$^{-1}$/H, whose element-11 solid-angle integral is
$C_{\rm sca,tot}-\eta\,C_{\rm sca,ref}$.
For the total extinction matrix the aligned population must not be counted twice:
the isotropic \code{Cext} from \code{dust\_extinction} already contains its
reference-orientation average $C_{\rm ext,ref}$, so assemble
\begin{equation}
  K_{\rm total} = (C_{\rm ext}^{\rm iso} - \eta\,C_{\rm ext,ref})\,I
                  + \eta\,K_{\rm al}(\theta_i),
\label{eq:ktotal}
\end{equation}
which returns $C_{\rm ext}^{\rm iso}$ on the diagonal at the incidence average and
adds the direction-resolved dichroism and birefringence. These aligned-scattering
quantities are in \um$^2$ (per H), whereas \code{dust\_extinction} returns cm$^2$/H;
convert with a factor $10^{8}$ before combining them, as
Eq.~(\ref{eq:ktotal}) requires.

\paragraph{Inlining.} The loaded grids and integrals are also exposed
\code{public, protected} (\code{scm\_lambda}, \code{scm\_theta\_i},
\code{scm\_theta\_s}, \code{scm\_phi}, \code{scm\_theta\_ran}, \code{scm\_cext\_tot},
\code{scm\_csca\_tot}, \code{scm\_cext\_ref}, \code{scm\_csca\_ref},
\code{scm\_F\_tot}, \code{scm\_F\_ref}, \code{scm\_Z}, and the \code{scm\_n*}
sizes), so a host that prefers to inline its own lookups can take them once at
startup and never call back. A complete two-cell reference consumer of every call
is \code{sed/rt\_example/use\_dustlib\_scatmat.f90} (\code{make
use\_dustlib\_scatmat.x}); it demonstrates the band selection,
Eq.~(\ref{eq:ktotal}), the aligned and random queries, and the closure checks.

\subsection{Division of labor}\label{sec:pol:division}

Table~\ref{tab:poldivision} is the part to get right. SEDust stops at the
intrinsic quantity; everything that depends on where the field points is left
to the transfer code, because SEDust knows nothing about the cell geometry.

\begin{table}[h]
\centering
\begin{tabular}{@{}p{0.46\textwidth}p{0.46\textwidth}@{}}
\toprule
Done by SEDust & Left to the calling code \\
\midrule
Integral over the grain size distribution &
$\sin^2\gamma$, with $\gamma$ the angle between the \emph{local magnetic
field} and the \emph{line of sight} (not a sky-plane angle) \\[2pt]
Weighting by the alignment efficiency $f_{\rm align}(a)$ &
Turbulent depolarization $F_{\rm turb}$ \\[2pt]
Summation over components (only astrodust is aligned; PAHs contribute zero) &
Splitting into $Q$ and $U$ using the position angle $\phi$ of the projected
field \\
\bottomrule
\end{tabular}
\caption{What is already integrated into \code{lamI\_pol} and into
\code{Cpol\_ext} (equivalently, column 8), and what a radiative transfer host
must still apply.}
\label{tab:poldivision}
\end{table}

\noindent The observed Stokes emissivities follow from the exported quantity as
\begin{equation}
(j_Q,\,j_U) \;=\; \code{lamI\_pol}\,\sin^2\!\gamma\;F_{\rm turb}\,
                  (\cos 2\phi,\;\sin 2\phi),
\label{eq:stokesuse}
\end{equation}
with $j_I = \code{lamI\_total}$ unchanged. The same $\sin^2\gamma$ and
$F_{\rm turb}$ multiply \code{Cpol\_ext} (column 8) when the extinction matrix
is assembled.
Equation~(\ref{eq:stokesuse}) is stated here as a usage rule; its derivation is
in \S5.1 of \code{aligned\_grain\_polarization.pdf}.

\subsection{Minimal example}

\begin{lstlisting}
use dust_lib
use radfield, only: J_Mathis
type(dust_model_t)    :: m
real(wp), allocatable :: J_lam(:), lamI_total(:), lamI_pol(:), p(:)
integer :: i

call build_astrodust(m, qtab, sizedist, 200, 2.7_wp, 5.0e3_wp)
allocate(J_lam(m%NLAM), lamI_total(m%NLAM), lamI_pol(m%NLAM), p(m%NLAM))

call J_Mathis(1.0_wp, m%lam, J_lam)
call dust_emission(m, J_lam, lamI_total, lamI_pol=lamI_pol)

! intrinsic polarization fraction (before any geometric factor)
do i = 1, m%NLAM
   if (lamI_total(i) > 0.0_wp) then
      p(i) = lamI_pol(i) / lamI_total(i)
   else
      p(i) = 0.0_wp
   end if
end do
\end{lstlisting}

\noindent Multiply \code{p} by $\sin^2\gamma\,F_{\rm turb}$ to obtain the
polarization fraction of a real line of sight.

\subsection{Limits}

\begin{itemize}
\item \textbf{Astrodust only.} \code{build\_dl07} and \code{build\_zubko} have
      no polarized optics. \code{build\_dl07} releases any polarized arrays left
      by a previous astrodust build, so \code{lamI\_pol} returns zeros rather
      than stale values.
\item \textbf{PAHs contribute zero} ($f_{\rm align}=0$, an HD23 assumption).
      The polarized emission therefore vanishes in the mid infrared, where
      stochastically heated PAHs dominate.
\item \textbf{Circular polarization: available only from the regenerated table.}
      The released optics table carries $Q_{\rm ext}$, $Q_{\rm abs}$ and
      $Q_{\rm sca}$ but no phase retardation, so with it $C_{\rm circ}=0$. The
      regenerated orientation-resolved table (\S\ref{sec:pol:source}) adds a
      birefringence block, and \code{dust\_extinction} returns it as the optional
      \code{Cbir\_ext} (\S\ref{sec:ext}); the aligned scattering table carries the
      same birefringence in its $K(\theta_i)$ (\S\ref{sec:pol:scatmat}).
      Consuming it in the $U\leftrightarrow V$ transfer term is the host's task.
\item \textbf{Scattering by aligned grains: now available.} The angle-resolved
      Mueller matrix of the aligned astrodust spheroid, which no released file
      contains, is computed by \code{run\_scatmat\_aligned.x} and served through
      the API of \S\ref{sec:pol:scatmat}. The release-derived cross sections still
      supply only the integrated $Q_{\rm sca}$; the angular optics come from the
      T-matrix engine.
\item \textbf{Far-infrared and submillimeter polarized emission is complete
      without a scattering matrix.} The albedo there is essentially zero, so the
      transfer needs only the extinction matrix and the emission vector, both
      fully determined by what is exported here.
\end{itemize}

\paragraph{Reference numbers.} For a Mathis ISRF the intrinsic polarization
fraction \code{lamI\_pol/lamI\_total} is $0.00\%$ at 12\um, $13.67\%$ at
100\um, $17.15\%$ at 154\um, $18.79\%$ at 350\um\ and $19.15\%$ at 850\um.
Column 8 reproduces the released \code{polarized\_extinction.dat} to a median
of $0.0296\%$.

%====================================================================
\section{The extreme-ultraviolet extension}\label{sec:euv}

The astrodust model's wavelength grid is the T-matrix $Q$ table's, which stops
at $0.0912$\um\ --- 13.6 eV, the Lyman limit. A host that transports ionizing
radiation needs shorter wavelengths, so both \code{build\_astrodust} and
\code{build\_dl07} take an optional \code{lam\_min} [\um]:
\begin{lstlisting}
call build_astrodust(m, qtab, sizedist, NT, T_lo, T_hi, lam_min=1.001e-4_wp)
\end{lstlisting}
Log-spaced points are then prepended from \code{lam\_min} up to the table, at a
step no coarser than the table's own $\Delta\ln\lambda = 0.01156$ at its
short-wavelength end, with \code{m\%lam(1)} set to \code{lam\_min} exactly.
\textbf{Omitting \code{lam\_min} prepends nothing and leaves the model
bit-identical to before.}

\subsection{Where the optics come from in the extended band}

\begin{itemize}
\item \textbf{Astrodust grains.} Bohren--Huffman Mie for the
      \emph{volume-equivalent sphere} on the Draine \& Hensley (2021)
      dielectric function --- the same material as the $Q$ table, but not the
      same shape. See \S\ref{sec:euv-shape}.
\item \textbf{PAH / carbonaceous.} Above 21.4 eV
      ($1/\lambda = 17.25\ \mu\mathrm{m}^{-1}$) the DL07 PAH cross section is
      zero by construction, so graphite is the \emph{whole} carbonaceous
      absorption there. Below that energy the D16 turbostratic-graphite sphere
      table is used as before; above it the code switches to Mie on the D03
      graphite dielectric functions (random orientation, $\tfrac13$ parallel
      $+\,\tfrac23$ perpendicular), because the D16 table's radius axis bottoms
      out at $10^{-3}$\um\ while the astrodust size grid starts at
      $3.55\times10^{-4}$\um. In the Rayleigh limit $Q_{\rm abs}\propto a$, so
      clamping a smaller grain to the table's first radius overestimates its
      absorption by $10^{-3}\,\mu\mathrm{m}/a$ --- up to $2.4\times$ for the
      smallest PAH the size distribution populates ($4.22\times10^{-4}$\um).
      \emph{This branch is unreachable on the unextended grid}: its shortest
      wavelength gives $1/\lambda = 10.96 < 17.25$, so every previous result is
      unaffected.
\item \textbf{DL07 model.} Its optics are Mie on the D03 dielectric functions at
      \emph{every} wavelength, so nothing changes character across the band and
      the extension is a grid extension only.
\item \textbf{Zubko model.} No extension, and none possible: its tabulated
      optics already span $10^{-3}$ to $10^{4}$\um, and those tables \emph{are}
      the model definition, so there is no dielectric function to extend them
      with.
\end{itemize}

\subsection{The shape approximation, and why it is bounded}\label{sec:euv-shape}

The $Q$ table is a random-orientation T-matrix solution for an oblate spheroid
of axial ratio $b/a = 1.4$; the extended band uses the volume-equivalent
\emph{sphere}. The justification is \textbf{not} the anomalous-diffraction
argument $|m-1|\ll1$: for this material $|m-1| = 0.765$ at the seam and still
$0.564$ at 20 eV. It rests instead on the two limits that actually apply:

\begin{itemize}
\item \textbf{Small grains are in the Rayleigh limit} ($x = 2\pi a/\lambda =
      0.033$ at the seam for the smallest populated radius, and only $0.24$ at
      100 eV). There the ellipsoid polarizability is
      $\alpha_j = V(\epsilon-1)/\{4\pi[1 + L_j(\epsilon-1)]\}$, so shape enters
      \emph{only} through $L_j(\epsilon-1)$. And $|\epsilon-1|$ falls from
      $1.86$ at 13.5 eV to $1.10$ at 20 eV, $0.28$ at 40 eV and $0.061$ at
      100 eV: the depolarization factors drop out and sphere and spheroid
      converge on the same $C_{\rm abs}$.
\item \textbf{Large grains are in the geometric-optics limit}, where $C_{\rm
      ext}\to2\times$(orientation-averaged projected area). By Cauchy's theorem
      that average is $S/4$, and for $b/a = 1.4$ the spheroid's $S/4$ exceeds
      $\pi a_{\rm eff}^2$ of its volume-equivalent sphere by 2.12\%. Mie is
      therefore low by $1 - 1/1.0212 = 2.08\%$ in that limit. \textbf{The error
      converges to a constant, not to zero}, however far into the EUV one goes.
\end{itemize}

\noindent Consequently the error does \emph{not} fall monotonically with photon
energy. Only $a\lesssim2\times10^{-3}$\um\ improves monotonically; $a\gtrsim
0.01$\um\ gets \emph{worse} between 13.6 and 30 eV before recovering, and the
largest sizes settle at a nearly energy-independent $-1.9$ to $-2.0\%$. Over the
whole band and the whole size range the error stays within about 2\%;
\code{sed/src/q\_astrodust.f90} tabulates it size by size.

\paragraph{Seam continuity.} In the size-integrated curve the step across
$0.0912$\um\ is $-1.54\%$ in $C_{\rm ext}$ and $-1.90\%$ in $C_{\rm abs}$ ---
\emph{smaller} than the neighboring steps of the table's own grid ($-1.95\%$ and
$-1.72\%$ for $C_{\rm ext}$; $-2.20\%$ and $-1.94\%$ for $C_{\rm abs}$), so
absorption and extinction join without a visible discontinuity. Scattering does
show one: $C_{\rm sca}$ steps by $+0.74\%$ where its neighbors step by $-0.39\%$
and $-0.36\%$, and the albedo by $+2.31\%$ where its neighbors give $+1.59\%$
and $+1.39\%$ --- a seam of roughly one percentage point, consistent with the
2\% shape error being carried mostly by $Q_{\rm sca}$.

\subsection{The polarized optics do not extend with the grid}\label{sec:euv-pol}

\textbf{This is the one place where the polarized branch has a hole, and it is
announced rather than hidden.} The dichroic and birefringent cross sections come
from the orientation-resolved table, which covers the $Q$ table's own
$0.0912$--$39810$\um\ and no further. When \code{lam\_min} carries the grid
below that, \code{build\_Cpol} looks for the EUV companion table
\code{q\_astrodust\_jori\_euv\_P0.20\_Fe0.00\_1.400.dat.gz}. \textbf{That table
does not ship}, so what normally happens is:
\begin{lstlisting}[style=sh]
build_Cpol: no EUV polarized table (...)
            dichroic extinction and birefringence are zero below 9.120E-02 um
            (zero by omission, not by physics).
\end{lstlisting}
The message goes to \code{error\_unit} unconditionally, and the EUV block of
\code{Cpol\_ext} and \code{Cbir\_ext} is left at zero. \textbf{A host must not
read that zero as a physical result.} Measured from first principles
(\code{tmatrix/driver/euv\_polarized\_optics.f90}), the alignment-weighted,
size-integrated $|C_{\rm pol,ext}|/C_{\rm ext}$ is $1.26\times10^{-3}$ at the
$0.0912$\um\ seam, \emph{rises} to $3.64\times10^{-3}$ near 20.6 eV --- a factor
2.9 --- and decays to $1.6\times10^{-4}$ at 100 eV, \textbf{with the sign
opposite to the optical band}; the reversal falls near $0.106$\um. A sphere has
exactly zero dichroism and exactly zero birefringence, so the scalar EUV optics
could not have supplied these entries even in principle.

To fill the band, generate the companion table with
\code{tmatrix/run\_q\_jori.x} (\S\ref{sec:standalone}) and pass the resulting
pair as \code{qpol\_euv\_path} and \code{qpol\_euv\_wave\_path}. With a table
present, \code{build\_Cpol} interpolates it in $\log\lambda$ and $\log a$, and a
grid running outside its $0.0124$--$0.0912$\um\ coverage is refused
(\code{status = 9}) rather than answered with an invalid limit. That table stops
at $0.0124$\um\ (100 eV) because the opaque geometric-optics limit it leans on
needs the chord optical depth $4\,{\rm Im}(m)\,x$ to be large, and for astrodust
that is $3.7$ at $x = 50$ there and falls below 1 by 200 eV.

The \emph{scalar} optics are unaffected by any of this: $C_{\rm ext}$,
$C_{\rm abs}$, $C_{\rm sca}$, \code{gbar} and \code{albedo} are complete over
the whole extended grid.

\subsection{The single-photon ceiling comes from the field}\label{sec:euv-ceiling}

The stochastic solvers set the top of a grain's enthalpy window from the energy
of the hardest single photon it can absorb. That ceiling used to be the
hard-coded Lyman limit, $13.6\,$eV, which an ionizing band makes wrong: a photon
harder than the ceiling drives an excursion the window cannot hold. It is now
\begin{center}
$\max\!\left(13.6\,\mathrm{eV},\ hc/\lambda_{\rm hard}\right)$, \qquad
$\lambda_{\rm hard} = \min\left\{\lambda_i\ :\ J_\lambda(\lambda_i) >
10^{-12}\,\max_j J_\lambda(\lambda_j)\right\}$,
\end{center}
computed by \code{hardest\_photon\_energy} in \code{sed/src/radfield.f90} and
called from all three places that need it --- \code{sed\_grain\_loop} and
\code{narrow\_iterative} in \code{sed\_astrodust.f90}, and
\code{qm\_solve\_grain} in \code{stoch\_qm.f90}. Any units may be passed for
$J_\lambda$, since only ratios within the array are used.

\paragraph{The ceiling is a property of the field, not of the grid.} That
distinction is the whole point, and a coincidence hides it. A model's wavelength
grid is the grid of its optics tables, which can reach far past the band a
transported field occupies:

\begin{center}
\begin{tabular}{@{}lll@{}}
\toprule
model & grid floor & implied photon energy \\
\midrule
astrodust & $0.0912$\um\ (T-matrix $Q$ table)        & $13.595\,$eV \\
DL07      & $0.0912$\um\ (the same table)            & $13.595\,$eV \\
Zubko     & $1.000\times10^{-3}$\um\ (DustEM tables) & $1239.8\,$eV \\
\bottomrule
\end{tabular}
\end{center}

\noindent For astrodust and DL07 the grid floor \emph{is} the Lyman limit, so
$13.595\,\mathrm{eV} < 13.6\,\mathrm{eV}$: the maximum returns the constant, and
grid and field give the same answer whichever one is measured. For Zubko they
differ by a factor of 91, and the field is the one that is right --- a Mathis
field carries no photon at all below $0.0912$\um, however far the optics table
extends.

\paragraph{Taking it from the grid costs resolution, not physics.}
\code{umax} sets the upper edge of a \emph{fixed} number of enthalpy bins, so a
ceiling 91 times too high starts the solve with bins 91 times too coarse. The
refinement loops do contract the window once the top bin is unpopulated, but the
contraction is guarded (\code{umax > 1.02*umaxlo} and
\code{umax > 1.01*ub(jcut)}) and capped at \code{MAX\_ITER = 10}, so it does not
return to where it began; expansion carries no such guard. Measured in a host
code on the Zubko model, with no \code{lam\_min} requested at all: dust
temperature $+0.0705\,$K median and positive in every cell, the ten brightest
SED bins $-0.9$ to $-1.9\%$, the emission image $1.9\%$ median and $5.4\%$
maximum, total emitted energy $-0.08\%$. A systematic redistribution at nearly
conserved energy is what a coarser bin grid produces, and
\code{docs/EUV\_EXTENSION\_HOST\_REGRESSION.md} records the measurement.

\paragraph{Why $10^{-12}$ of the peak, and not simply $J_\lambda>0$.} Exact
zeros must be excluded --- \code{J\_Mathis} returns exactly zero below
$0.0912$\um\ --- but an exact positivity test is too fragile: a host that hands
over denormal or roundoff-level residue in its unilluminated bins would push the
ceiling back up by orders of magnitude. A component a fraction $10^{-12}$ below
the peak of $J_\lambda$ contributes at most that fraction of the photon
absorption rate; allowing two decades for the run of $C_{\rm abs}$ across the
illuminated band, the enthalpy states it could populate carry probability
$\lesssim10^{-10}$ of the peak, at or below the $10^{-13}$ tail level
(\code{PMIN\_UP}, \code{PMIN\_UP\_QM}) at which both solvers already truncate
the window. Such a component cannot change a resolved excursion, and $10^{-12}$
still sits eighteen decades above the residue it is meant to reject.

\medskip
\noindent The error is one-sided on purpose. Both refinement loops \emph{expand} the
window while the top bin is still populated, so an underestimated ceiling is
corrected as the loop runs; the guarded contraction can stop short of where it
began, so an overestimate is not corrected at all. \code{MAX\_ITER = 10} bounds
the correction in either direction. Overestimating is therefore the dangerous
direction, which is why the threshold is set high enough to reject junk rather
than as low as floating point allows.

\paragraph{\code{main\_zubko.x} exists to keep this path covered.} The defect
escaped \code{main\_astrodust.x} and \code{main\_dl07.x} because for those two
models the grid floor and the field's short end are the same number: no driver
in the tree ran the emission solvers on a model whose optics grid reaches past
the illuminated band. \code{main\_zubko.x} (\S\ref{sec:standalone:zubko}) closes
that gap and is built by the default \code{make}. It prints the two candidates
side by side, so the difference shows in the first lines of output:
\begin{lstlisting}[style=sh]
$ ./main_zubko.x            # Mathis field, stops at the Lyman limit
 optics grid short end :  1.0000E-03 um  ->   1.23984E+03 eV
 hardest photon in the field:    1.35946E+01 eV
 single-photon bound in use :    1.36000E+01 eV
$ ./main_zubko.x euv        # a hard component added below the Lyman limit
 optics grid short end :  1.0000E-03 um  ->   1.23984E+03 eV
 hardest photon in the field:    3.70141E+02 eV
 single-photon bound in use :    3.70141E+02 eV
\end{lstlisting}
Without \code{euv} the ceiling must stay at $13.6\,$eV however far the optics
grid extends, and the Zubko SED is then identical to what the hard-coded
constant produced; with \code{euv} it must follow the field up, and it does.

\paragraph{A model built without \code{lam\_min} is unaffected.} For astrodust
and DL07 the maximum selects the $13.6\,$eV constant on the unextended grid
either way, so those results are unchanged to the last bit.

\subsection{What else the extension does, and does not, change}

\begin{itemize}
\item \textbf{The Planck integrals are anchored to the table range.} The
      internal $\log\lambda$ integration grid for $\kappa_B(T,a)$ keeps its step
      and its sample points over the optics-table range and merely
      \emph{prepends} points below it, so widening the model grid cannot
      coarsen the sampling of the range that carries the integrand. Measured:
      with a $1.001\times10^{-4}$\um\ floor, $\kappa_B$ is unchanged to
      round-off over the thermal range ($\max|\Delta\kappa_B|/\kappa_B =
      3\times10^{-16}$ for $T\le3000$ K, and exactly zero below $\sim2400$ K);
      $\kappa_{\rm CMB}$ is bit-identical. The only nonzero difference is
      $1.1\times10^{-8}$ at the very top of a 5000 K grid, where the Wien tail
      genuinely reaches into the extended band --- that is real extra emission,
      not a sampling artifact. Without the anchoring, spending a fixed budget of
      points on the widened interval would have moved $\kappa_B$ by
      $\sim10^{-4}$ for a purely numerical reason.
\item \textbf{Stochastic heating by hard photons is not solved by this change.}
      The extension fixes the photon-energy mapping and the absorbed-energy
      accounting. It does not address the finite enthalpy grid: upward
      transitions above the top of that grid are accumulated into the highest
      temperature bin, and for the smallest PAHs a 54 eV photon already exceeds
      $H(5000\ {\rm K})$. Nor does it add photoionization, electron energy loss,
      fragmentation or PAH destruction. Treat the extended band as extinction
      and absorbed-energy physics, not as a validated hard-EUV PAH
      \emph{emission} model.
\end{itemize}

\subsection{How far each model can be checked}\label{sec:euv-refs}

\begin{table}[h]
\centering
\small
\begin{tabular}{@{}L{0.30\textwidth}L{0.24\textwidth}L{0.38\textwidth}@{}}
\toprule
author-provided data & coverage & what it contains \\
\midrule
Draine, \code{kext\_albedo\_WD\_MW\_3.1\_60\_D03} (DL07 / WD01) &
$10^{-4}$--$10^{4}$\um\ $=1.24\times10^{-4}$\,eV--12.4\,keV &
$\lambda$, albedo, $\langle\cos\theta\rangle$, $C_{\rm ext}$/H, $K_{\rm abs}$,
$\langle\cos^2\theta\rangle$ \\
HD23 astrodust release (\code{extinction.dat}, \code{scattering.dat},
\code{wav.dat}) &
$0.1$--$3\times10^{4}$\um\ $=4.1\times10^{-5}$--\textbf{12.4\,eV} &
absorption and scattering $\tau/N_{\rm H}$ only. \textbf{No
$\langle\cos\theta\rangle$ product exists anywhere in the release.} \\
Zubko ZDA optics tables &
$10^{-3}$--$10^{4}$\um\ $=1.24\times10^{-4}$\,eV--\textbf{1.24\,keV} &
$Q_{\rm abs}$, $Q_{\rm sca}$, $Q_{\rm ext}$, $g$ for each radius \\
\bottomrule
\end{tabular}
\caption{What the model authors published, and how far it reaches.}
\label{tab:refcoverage}
\end{table}

\noindent The consequence is blunt: \textbf{the astrodust EUV band has no
external reference at all}, because the HD23 release stops at 12.4 eV --- exactly
where the extension begins. Its accuracy claim rests entirely on the internal
argument of \S\ref{sec:euv-shape} and on the seam being continuous.

The DL07 model \emph{can} be checked, because Draine published the same model
over the whole range. Comparing \code{data/kext\_dl07\_MW\_euv.dat} with
\code{kext\_albedo\_WD\_MW\_3.1\_60\_D03.all\_2003} on the reference's own
wavelengths (and counting only points where the reference grid is no finer than
ours --- elsewhere the comparison measures our grid, not our optics) gives, per
decade in $\lambda$, mean $|\Delta C_{\rm ext}|/C_{\rm ext}$ of
0.056--0.86\%, mean $|\Delta\,\mathrm{albedo}|\le0.0019$ and mean
$|\Delta\langle\cos\theta\rangle|\le0.00054$. \code{calc\_kext.x dl07 euv}
prints that table. The two residual features are worth naming: the excluded
points cluster at the X-ray absorption edges (Fe\,K, Si\,K, Mg\,K, Fe\,L, O\,K,
C\,K), which the reference brackets far more finely than our
$\Delta\ln\lambda = 0.0116$ log grid can resolve; and the largest single
deviation, 7.9\% in the 1--10\um\ decade, is confined to the 6.2, 7.6 and
7.9\um\ PAH C--C features, where Draine's 2003 table uses a different vintage of
the PAH cross sections.

%====================================================================
\section{Standalone programs}\label{sec:standalone}

\code{make} in \code{SEDust/sed/} builds five executables. All of them are run
\emph{from} \code{SEDust/sed/}, because every data path is relative to
\code{../}.

\subsection{\code{calc\_kext.x} --- size-integrated transport optics}

Builds a model and calls \code{size\_integrated\_extinction} on it, then writes
the result under \code{../data/}. Every model goes through the same two calls,
so the files it writes \emph{are} the optics an RT host gets in memory --- these
are the files \code{dust\_extinction} reads back (\S\ref{sec:ext}).

\begin{lstlisting}[style=sh]
./calc_kext.x astrodust [euv | <lam_min in um>]
./calc_kext.x dl07      [euv | <lam_min in um>]
./calc_kext.x zubko     [formula | table]
./calc_kext.x from_files <descriptor> [data_dir]
\end{lstlisting}

\noindent With no argument it prints its usage and stops. \code{euv} carries the
grid down to whatever the model's own dielectric function covers; an explicit
number requests a specific floor in \um. \code{zubko} defaults to
\code{formula}. \code{from\_files} takes the descriptor path, and the data
directory defaults to the descriptor's own directory.

Columns are
\begin{center}
\code{lambda[um] \ albedo \ <cos> \ C\_ext/H \ C\_abs/H \ C\_sca/H}
\end{center}
in \um\ and cm$^2$ per H nucleon, followed by $K_{\rm abs}$ [cm$^2$/g] and ---
for \code{kext\_astrodust\_MW.dat} only --- an eighth column $C_{\rm polext}$/H
(Table~\ref{tab:kextcols}). The first four columns are in the order of Draine's
\code{kext\_albedo} tables, so a reader written for those files (for instance
the \code{par\%ion\_dust\_kext} reader of the MoCHII host, which takes four
reals from each row and uses $\lambda$, albedo and $C_{\rm ext}$) parses these
unchanged. Note the fifth column differs from Draine's: his is $K_{\rm abs}$
[cm$^2$/g], ours is $C_{\rm abs}$/H [cm$^2$/H], so no dust-mass normalization
is folded into the transport columns.

\paragraph{The $K_{\rm abs}$ column.} Every product carries it, for every model.
It is $C_{\rm abs}/{\rm H}$ divided by the model's dust mass per H,
Eq.~(\ref{eq:massperh}), which the model reports through
\code{dust\_mass\_per\_H} --- so the column and the library agree by
construction rather than by a constant copied between them. The normalization is
one wavelength-independent number, as well defined in the extreme ultraviolet as
in the optical, and the header of each file states its value and where that
model's grain densities came from. The one product that would not get the column
is a file-defined model whose optics tables declare no density; then
$M_{\rm dust}/N_{\rm H} = 0$, the program says so, and writes the six transport
columns alone.

\paragraph{Why the EUV products carry no polarized column.} Of the four
products, only \code{kext\_astrodust\_MW.dat} --- the unextended one --- has the
dichroic column. Below $0.0912$\um\ that quantity is the announced zero of
\S\ref{sec:euv-pol}, and writing it into a data file would turn a documented
deficit into a number a host cannot tell from a measurement. The EUV products
are therefore scalar transport optics only, and are byte-for-byte the same files
the scalar branch (v1.00) writes --- which is also the cleanest demonstration
that the scalar path is identical in the two branches.

After writing, the program checks $C_{\rm ext} = C_{\rm abs}+C_{\rm sca}$ and
the ranges of albedo and $\langle\cos\theta\rangle$, and --- for astrodust and
DL07 --- prints the comparison against the author-provided data of
Table~\ref{tab:refcoverage}.

\paragraph{These files are transport optics, not a dust model.} They carry no
size-resolved $C_{\rm abs}(\lambda,a)$, no grain enthalpy and no Planck
integral, so stochastic heating and thermal emission cannot be computed from
them. A host that needs emission must build the model and take the same curves
from \code{dust\_extinction}, which is what keeps the absorbed power and the
re-emitted spectrum referring to one and the same grain.

\subsection{\code{main\_astrodust.x} --- the astrodust production SED}

Solves the astrodust+PAH SED for a single Mathis-ISRF cell at $U = 1.585$
($\log U = 0.20$, the HD23 best fit), on a 200-point log $T$ grid from 2.7 to
5000 K. Reads the T-matrix $Q$ table
\code{../tmatrix/output/q\_astrodust\_P0.20\_Fe0.00\_1.400.dat} and
\code{../data/release/size\_distribution.dat}; writes
\code{output/astrodust\_irem\_ours\_<tag>S1.dat}, \code{\ldots S2.dat} and
\code{\ldots PAH.dat}, each with $\lambda$ [\um] and \lamIlam/\Nh.
Optional arguments, in any order, each of which tags the output filename so a
non-default run cannot clobber the production one:
\code{qm}, \code{qm\_dbcon}, \code{draine} (solver);
\code{logU=X}; \code{nstate=N}, \code{nisrf=N} (QM solver sizes);
\code{induced} (turn the stimulated-emission factor on); \code{photcut};
\code{c2} (Stage-1 density-corrected prefactor); \code{mathis\_orig}.

The tag is assembled after all the arguments have been read, in one fixed
order --- solver, enthalpy stage, radiation field, emission term, numerical
settings --- so the order they are typed in does not matter:
\code{induced qm} and \code{qm induced} both write
\code{\ldots\_qm\_induced\_S1.dat}. The three solver options select the same
thing, so naming two of them is refused rather than resolved by taking the
last. The driver prints the filename it will write before it starts solving.

\subsection{\code{main\_dl07.x} --- the DL07 reference SED}

Solves the Draine \& Li (2007) silicate + carbonaceous SED at $U = 1$ with the
WD01 size distribution. \code{./main\_dl07.x [model] [qm|draine]}, where
\code{model} is one of \code{mw31\_00 \ldots mw31\_60} (default
\code{mw31\_60}), \code{lmc2\_00}, \code{lmc2\_05}, \code{lmc2\_10},
\code{smc}; an unknown name is rejected with the valid list. Writes
\code{output/dl07\_sed\_ours\_<model>[\_qm|\_draine].dat} with columns
$\lambda$, total, silicate and carbonaceous \lamIlam/\Nh.

\subsection{\code{main\_zubko.x} --- the Zubko/ZDA SED, and the wide-grid
regression}\label{sec:standalone:zubko}

Solves the Zubko, Dwek \& Arendt (2004) BARE-GR-S SED --- PAH + graphite +
silicate, the ZDA size-distribution formula, the DustEM optics and calorimetry
tables --- for a Mathis ISRF at $U = 1$, and writes
\code{output/zubko\_sed\_ours[\_<solver>][\_euv].dat} with columns $\lambda$,
total, PAH, graphite and silicate \lamIlam/\Nh.
\begin{lstlisting}[style=sh]
./main_zubko.x [heuristic | draine | qm | equil] [euv]
\end{lstlisting}
\noindent The solver defaults to \code{heuristic}, the model default.
\code{euv} replaces the Mathis field below the Lyman limit by a diluted
$10^5\,$K blackbody, so that the field really does carry photons above
$13.6\,$eV; the dilution is chosen so that band deposits roughly as much power
as the whole Mathis field --- enough for the hard photons to matter, little
enough that the small grains still heat stochastically.

\medskip
\noindent This is the emission counterpart of \code{./calc\_kext.x zubko}, which exercises
only the extinction size integral. It exists because Zubko is the one shipped
model whose optics grid ($10^{-3}$\um) reaches far past the band a transported
field occupies ($0.0912$\um\ for the Mathis ISRF), so anything in the
stochastic-heating path that reads the grid where it should read the field shows
up here and nowhere else. Running it with and without \code{euv} is the
regression of \S\ref{sec:euv-ceiling}.

\subsection{\code{make\_enthalpy.x} --- enthalpy tables (diagnostic)}

Tabulates $U(T, a_{\rm eff})$ for the two astrodust enthalpy stages and writes
\code{output/enthalpy\_S1.dat} and \code{output/enthalpy\_S2.dat} (201
log-spaced temperatures from 2.7 to 5000 K $\times$ the 169 radii of
\code{../data/dielectric/DH21\_aeff}, $T$ outer and $a$ inner). It takes no
arguments. \textbf{It is not required to compute an SED}: the solver calls the
closed-form Debye sums directly in its inner loop. Use it to compare the two
stages, to plot $U(T)$, or to spot-check against published DL01 values.

\subsection{Polarized and diagnostic programs}

Four more targets are built on request rather than by the default \code{make}:
\code{calc\_polext.x} (polarized extinction per H, checked against the HD23
release), and the three self-checking test programs
\code{test\_scalar\_mode.x}, \code{test\_scatmat\_aligned.x} and
\code{test\_euv\_extension.x}. The tests print \code{ALL CHECKS PASSED} or the
failing check; build them explicitly, since \code{make} alone will not refresh a
stale binary:
\begin{lstlisting}[style=sh]
make test_scalar_mode.x test_scatmat_aligned.x test_euv_extension.x
\end{lstlisting}

\subsection{Polarized optics at other wavelengths}\label{sec:standalone:qjori}

These programs live in \code{SEDust/tmatrix/}, not \code{sed/}, and are run from
there. They exist because \textbf{the library reads tables and never runs the
T-matrix at run time} --- not to keep the code small, but because convergence
cannot be certified above $x\simeq55$ (below), and a silently wrong value is
worse than an absent one. A revision of the T-matrix core is planned, so read
this as the current policy rather than a permanent one. The core has since been
rewritten as a re-entrant modern-Fortran library, but that change was structural
and its output byte-identical; the convergence ceiling it was about is untouched.

Polarized \emph{transfer} at every wavelength is rarely wanted; the five optical
bands that ship cover the usual case. What matters is that any other wavelength
can be computed:
\begin{lstlisting}[style=sh]
./run_q_jori.x lam L1 [L2 ...] [ja=JA1:JA2] [tag=NAME]
./run_q_jori.x lamfile PATH    [ja=JA1:JA2] [tag=NAME]   # one per line, '#' comments
./run_q_jori.x lammerge STEM FILE [FILE ...]             # reassemble the ja= windows
\end{lstlisting}
Each writes a pair,
\code{q\_astrodust\_jori\_P0.20\_Fe0.00\_1.400.TAG.dat} and \code{.TAG.wave},
which is exactly what \code{qpol\_euv\_path} and \code{qpol\_euv\_wave\_path}
take. Because \code{build\_Cpol} interpolates in $\log\lambda$, the pair needs
\textbf{at least two wavelengths}.

\code{ja=JA1:JA2} splits the 169 radii across \emph{processes}, not threads. That
was once an engine constraint: the f77 solver handed the converged T-matrix to
\code{AMPL} through \code{COMMON /TMAT/} and kept further working storage in
COMMON blocks, two of them blank, so the core was not re-entrant. None of that
state is left --- every working array is in a caller-owned
\code{tmatrix\_workspace\_t}, one active caller for each workspace, and distinct
workspaces may be evaluated concurrently --- but \code{run\_q\_jori.x} keeps the
process windows, which are already measured. Measured on this machine: 9m22s of
wall time for 34m48s of processor time, 57 processes over three wavelengths.

\paragraph{Convergence is certified, not assumed.} In \code{lam},
\code{lamfile} and \code{euv} mode every node with $x\in(50,60]$ is solved twice
--- $(\code{DDELT},\code{NDGS}) = (10^{-3},2)$ and $(3\times10^{-3},3)$ --- and
rejected if the two disagree by more than 5\%, in which case the
geometric-optics limit is used instead. Above $x = 60$ that limit is used
outright. This matters because \textbf{\code{IERR = 0} is not a sufficient
condition for convergence above $x\simeq55$}: at $\lambda = 0.0912$\um,
$a = 0.9441$\um\ ($x = 65.04$) the solver returns \code{IERR = 0} together with a
value about 50\% away from its neighbors. At the seam $x = 51.7$ and $54.7$ pass
the cross-check while $57.9$, $61.4$, $65.0$ and $68.9$ fail. The older
full-sweep and \code{range} modes keep the previous $x > 50$ rule, so they still
reproduce the shipped table byte for byte.

\paragraph{What certification costs.} Leaving $x>60$ at the geometric-optics
zero discards part of the dichroic extinction: nothing at the seam, $\sim8\%$ of
the band total at $0.031$\um, $\sim20\%$ at $0.022$\um\ and $\sim46\%$ at
$0.0124$\um. The true value lies between the certified $1.06\times10^{-4}$ and
the $1/x$ extrapolation $3.14\times10^{-4}$ --- already about 1\% of the
$5.5\times10^{-2}$ the V band reaches, so the residual is a small fraction of an
already small number. $C_{\rm pol}$, which drives polarized \emph{emission}, has
no such gap: the geometric-optics limit obtains it from the opaque-grain Fresnel
surface integral.

\paragraph{Scattering matrices are less flexible.}
\code{run\_scatmat\_aligned.x} already accepts wavelengths on the command line
(\code{./run\_scatmat\_aligned.x 0.44 0.55 0.79}), but three things
\code{run\_q\_jori.x} has are missing: its output stem is fixed, so a custom run
\textbf{overwrites} the shipped five-band table; there is no merge mode, so
bands cannot be \emph{added} to an existing table (a full recomputation costs
about 4.5 minutes per band); and it has neither a radius split nor the $x>50$
cross-check above. Adding a band today therefore means regenerating the whole
table and accepting the uncertified large-$x$ treatment.

%====================================================================
\section{The model object and accessors}

\begin{lstlisting}
type :: dust_model_t
   character(len=32)              :: name
   integer                       :: NA, NLAM, NT
   real(wp), allocatable         :: lam(:)        ! (NLAM) [um]
   real(wp), allocatable         :: T_first(:)    ! (NT)   [K]
   type(grain_pop_t), allocatable :: pops(:)
   integer                       :: n_channel
   character(len=16), allocatable :: channel_name(:)
   logical                       :: use_induced_emission   ! default .false.
   character(len=16)             :: stoch_method           ! default 'heuristic'
   logical                       :: verbose                ! default .false.
end type
\end{lstlisting}

\paragraph{Diagnostics.} \code{m\%verbose} (default \code{.false.}) keeps the
library solve path silent; set it \code{.true.} to print the same solver
diagnostics as the CLI drivers.

Convenience accessors (all in \code{dust\_lib}):
\begin{lstlisting}
n  = dust_nlam(m)              ! number of wavelengths
nc = dust_n_channel(m)         ! number of output channels
lam = dust_lambda(m)           ! allocatable copy of the wavelength grid [um]
nm  = dust_channel_name(m, ic) ! name of channel ic
md  = dust_mass_per_H(m)       ! dust mass per H nucleon [g/H]
\end{lstlisting}

\noindent \code{dust\_mass\_per\_H} is the conversion from the library's
per-H cross sections to a mass opacity, Eq.~(\ref{eq:massperh}) and
\S\ref{sec:ext}. Each population carries the solid density of its material
in \code{m\%pops(ip)\%rho\_bulk} [g\,cm$^{-3}$]; a population left at $0$
contributes nothing to the sum.

%====================================================================
\section{The generic file loader}

\code{build\_from\_files} reads a small text \emph{descriptor} and assembles a
model from data files, with no code changes. One \code{pop:} line per
population:
\begin{lstlisting}[style=sh]
# descriptor.txt
name = mymodel
# pop: <grain_type> <channel> <optics> <dnda_table> <calorimetry> <rho(0=use file)>
pop: pah 1 PAH_28_1201_neu.dat  SzDist_PAH.dat Graphitic_Calorimetry.dat 0.0
pop: gra 2 Gra_121_1201.dat     SzDist_Gra.dat Graphitic_Calorimetry.dat 0.0
pop: sil 3 suvSil_121_1201.dat  SzDist_Sil.dat Silicate_Calorimetry.dat  0.0
\end{lstlisting}
\begin{itemize}
\item \code{grain\_type}: \code{'sil'}, \code{'pah'}, or \code{'gra'} (selects the
      heat-capacity/mode model used by the \code{'qm'} solver; ignored by the GD
      solvers, which use the supplied enthalpy directly).
\item \code{channel}: integer output channel ($1\dots$); populations sharing a
      channel are summed.
\item \code{optics}: a DustEM/Zubko $Q$-table ($Q_{\rm abs}$ per radius and
      wavelength); $C_{\rm abs}=Q_{\rm abs}\pi a^2$.
\item \code{dnda\_table}: a 2-column table $a$ [\um] and ${\rm d}n/{\rm d}a$
      [cm$^{-1}$ H$^{-1}$].
\item \code{calorimetry}: a specific-enthalpy table $u(T)$ [erg/g];
      $H(T,a)=u(T)\,\rho\,\tfrac{4\pi}{3}a^3$.
\item \code{rho}: grain density [g\,cm$^{-3}$]; \code{0.0} uses the value in the
      optics-file header. The same number sets the population's
      \code{rho\_bulk}, so it fixes both the enthalpy above and the model's dust
      mass per H, Eq.~(\ref{eq:massperh}); an optics file that declares no
      density and a descriptor that supplies none leave the population out of
      that mass.
\end{itemize}
All files are sought under \code{data\_dir}. The coded builders
(\code{build\_astrodust}/\code{dl07}/\code{zubko}) are the recommended path for
the built-in models; \code{build\_from\_files} is for adding new data-defined
models.

%====================================================================
\section{Linking into a 3D RT code}

Compile your RT source against \code{libsedust.a} with the \code{.mod} search path
pointing at \code{SEDust/sed/}:
\begin{lstlisting}[style=sh]
gfortran -I<sed> my_rt.f90 <sed>/libsedust.a -fopenmp -o my_rt.x
\end{lstlisting}
Three complete examples are provided under \code{SEDust/sed/rt\_example/}, each
with its build command in its own header. No \code{ISO\_C\_BINDING} is needed
--- the RT code simply \code{use}s the Fortran module \code{dust\_lib}.
\begin{itemize}
\item \code{use\_dustlib.f90} --- the scalar host flow, and the one to read
      first: build a model once (which is also where its extinction table is
      attached), take the transport optics from
      \code{dust\_extinction} (including \code{gbar} and \code{albedo}), take
      one cell's emission from \code{dust\_emission}, then build a second model
      with \code{lam\_min} to show how a host that transports ionizing
      radiation widens the grid. It is byte-for-byte the same file in the
      scalar branch (v1.00).
\item \code{use\_dustlib\_pol.f90} --- where SEDust stops and the RT starts:
      \code{lamI\_pol} and \code{Cpol\_ext}, and the $\sin^2\gamma$, $F_{\rm
      turb}$ and position-angle factors the host must supply
      (\S\ref{sec:pol:division}).
\item \code{use\_dustlib\_scatmat.f90} --- the aligned-scattering query API
      along a two-cell photon path (\S\ref{sec:pol:scatmat}).
\end{itemize}

\paragraph{Threading.} \code{m} is read-only during \code{dust\_emission}; all
scratch is local. The solver uses OpenMP internally over grains. If your RT
parallelizes over cells, take care to avoid oversubscription (nested
parallelism); the default \code{'heuristic'} solver is serial per call, which
suits an outer cell-level parallel loop.

%====================================================================
\section{Units and conventions}

\begin{itemize}
\item Wavelengths \code{lam} are in \um\ throughout; cross sections in cm$^2$.
\item \code{J\_lam} is a mean intensity with the same units as the Planck
      function $B_\lambda$ (specific intensity per steradian). The library's
      output \lamIlam$/\Nh$ is then in
      erg\,s$^{-1}$\,cm$^{-2}$\,sr$^{-1}$\,H$^{-1}$.
\item The CMB (2.725\,K) is included in the heating field by \code{J\_Mathis}.
\item The induced-/stimulated-emission factor is off by default (the output is
      \emph{net} emission, matching the HD23 and DL07spec released SEDs).
\end{itemize}

%====================================================================
\section{Data files}

Model data lives under \code{SEDust/data/}:
\begin{itemize}
\item astrodust / DL07: the T-matrix $Q$-table
      (\code{../tmatrix/output/q\_astrodust\_*.dat}) and the HD23 release size
      distribution (\code{../data/release/size\_distribution.dat}).
\item Zubko: \code{../data/zubko/} --- the ZDA config/formula, the tabulated
      size distributions, the DustEM $Q$-tables, and the calorimetry tables
      from the Camps et al.\ 2015 benchmark; a ready descriptor is
      \code{../data/zubko/zubko\_descriptor.txt}.
\item Polarized optics: \code{../data/dielectric/q\_DH21Ad\_*.dat.gz} (the HD23
      release orientation-resolved $Q$ table, the default) and
      \code{../tmatrix/output/q\_astrodust\_jori\_*.dat.gz} (SEDust's own
      regenerated one, opt-in through \code{qpol\_path}, and the only one with a
      birefringence block); plus
      \code{../tmatrix/output/scatmat\_aligned\_*.dat.gz} (aligned scattering
      matrix), with their grid axes \code{../data/dielectric/DH21\_wave} and
      \code{DH21\_aeff}. The EUV companion table
      \code{q\_astrodust\_jori\_euv\_*.dat.gz} does \emph{not} ship
      (\S\ref{sec:euv-pol}); \code{DH21\_wave\_euv} is the axis it would be
      computed on.
\item Size-integrated transport optics, written by \code{calc\_kext.x}, see
      \S\ref{sec:standalone} --- \code{../data/kext\_astrodust\_MW.dat} (1129
      rows, $0.0912$--$39810$\um, with the dichroic 8th column),
      \code{kext\_astrodust\_MW\_euv.dat} (1719 rows, from
      $1.001\times10^{-4}$\um), \code{kext\_dl07\_MW.dat} (1129 rows,
      $0.0912$--$39810$\um), \code{kext\_dl07\_MW\_euv.dat} (1761 rows, from
      $6.205\times10^{-5}$\um) and \code{kext\_zubko\_BARE\_GR\_S.dat} (1201
      rows, $10^{-3}$--$10^{4}$\um). A host that reads a file reads one of
      these; a host that links the library gets the same curves from
      \code{dust\_extinction}, which reads one of these files too --- the two
      EUV products and the Zubko one are the defaults, see \S\ref{sec:ext}.
      The DL07 and astrodust unextended products carry exactly the
      long-wavelength rows of their EUV counterparts, bit for bit.
\end{itemize}

\end{document}
